\gddchapter{Interview sequence}


Below is the guideline for the semi-open guided interview. 
The interview starts with easy warm-up questions and progresses towards harder questions as it goes along. 

The answers were translated from polish back to english.

\section{Introduction}
Remind the participant that this is not a test, rather a joint exploration.

\section{Icebreakers}

\begin{QandA}
   \item Now that you’ve tried both versions, how are you feeling? Was there anything that immediately surprised you?
         \begin{answered}
            I mean the fact that in the one you can control with your keyboard you always know at first what you can do.
            Everything is predictable.
            There's almost no need to read the text or anything.
            You could just speedrun through the buttons and see what happens.
            In the speech one you cannot do that because the only way to know what you can do is from the text.
            I thought for a second that maybe I could, for example, take some other items than in the previous version.
            Because they were present in the text, I could assume, for example, in the shed, there is no way of knowing that the crowbar is the thing you should get if you don't go to the entrance first, if you know what I mean.
            If I played it the other way around, I don't think I could find the brother in the speaking one.
            besides the bugs that happened, just like from the gameplay, in the 10 minutes that it took me to get him in and out, I don't think I could do that in the speech one.
            \end{answered}

   \item If you had to describe the difference between these two games to a person who hasn't played either, what would you say?
         \begin{answered}
             I mean, I would describe them both as like a text adventure kind of game.
            The first one I would describe as pretty simple because of the way the game says what you can do.
            The other one is a bit more complex and you need to think about solutions a bit more.
            So it's more interesting in that sense for me, and I think that's how I would describe it to someone.
         \end{answered}
\end{QandA}




\section{Grand tour questions}
(remember to pause here 3-5 seconds)

\begin{QandA}


   \item From the moment you first used your voice to navigate, what were you thinking as you decided what to say to the system?

         \begin{answered}
            I thought it should be as easy to understand, as simple as it can be.
            To not use like complex sentences and things like that, just go say the name of the room, the name of the thing and like a simple command.
            Gabriel: And how was it like to use the voice input in the game?
            I mean, it was fine.
            I have not ever, I think, used the voice like strictly to input things into a game.
            So it was something different.
            But yeah, in this formula of a game, it could be fun to use.
         \end{answered}

    \item How was it like to use the voice input in the game?
         \begin{answered}
            I mean, it was fine.
            I have not ever, I think, used the voice like strictly to input things into a game.
            So it was something different.
            But yeah, in this formula of a game, it could be fun to use.
         \end{answered}


\end{QandA}




\section{Core comparative questions}

\begin{QandA}
   \item In which version did you feel more 'inside' the story world?
         \begin{answered}
            I think in the talking one, because you had to immerse more in the text and more focus on where you are, what you're doing.
            When you talk to the computer you feel like you have a conversation with the game, so it's also pretty interesting.
            But you need to be more focused, you cannot rely on the game telling you what to do.
         \end{answered}

   \item How did the 'effort' of thinking of what to say naturally compare to the 'effort' of knowing which keys to press?

         \begin{answered}
            I mean, both were pretty simple.
            As I said, I tried to make the commands as simple and on point as possible.
            So for me, it wasn't much difference.
            like just input wise because it was a simple very simple things freedom to say anything make you feel more in control
         \end{answered}

    \item Did the freedom to say 'anything' make you feel more in control, or was it more overwhelming than having a fixed list of keys?

         \begin{answered}
            I mean, it wasn't overwhelming for me.
            It was more like a freedom that I could think of another solution.
            I could try something else.
            I really like in like other video games when you can, for example, pick some random item at the beginning of the game and then two hours later it turns out you can do something with it.
            It's pretty cool and satisfying feeling.
            So when I don't have a strict list of things I can do, I can have that feeling that I can do something with it and try to find other solutions.
            So it's pretty fun.
         \end{answered}
    \item How would you compare the overall experience of using the two game versions? Which one would you choose and why?

         \begin{answered}
            I mean, how I said, the first one, the one with the keyboard inputs is pretty simple and you could go through it without really thinking about it, if you want.
            But the second one is like, you need to be taken in the story to be able to play it, because without reading and being invested somehow in this world, you cannot even move forward.
            So if I would want to play one of those games, I would want to play the second one, the voice one, I think.
         \end{answered}
\end{QandA}






\section{Probing}
In this part of the interview look into the sticky moments noted above and ask follow-up questions probing for more detail.
Include lexical reflection (eg. you used the word "clunky" when explaining navigation. Could you tell me more about what was happening
in that moment?)


\section{Final reflection}

\begin{QandA}
   \item During the play-through, there were a lot of, I mean, I know that four separate instances of, you know, you saying something and the game doing something completely different, right? How did that make you feel in those moments?
         \begin{answered}
            I mean, a bit frustrated, as every time something bugs and doesn't work as it should be.
            Also, it was a bit funny that the game, I said something, the game did something completely different.
            maybe for a moment I thought about it like it could it could be intended as such like have fun with the player to like invert expectations a bit it's like a breaking of a fourth wall kind of thing but maybe it's just my imagination running wild next time I have to explain myself to a committee
         \end{answered}

   \item After reflecting on everything today, is there anything else you’d like to add about using your voice to play this game that we haven't talked about?

         \begin{answered}
            uh i don't know for for this moment i don't think so
         \end{answered}
\end{QandA}

\section{Closing}
Thank the participant for his time and tell them again how much value his input and time gives me.






